\documentclass[11pt,a4paper]{article}
\usepackage[utf8]{inputenc}
\usepackage[T1]{fontenc}
\usepackage{geometry}
\geometry{margin=1in}
\usepackage{booktabs}
\usepackage{hyperref}
\usepackage{tabularx}
\title{t27b: A Fast Baseline AArch64 Compiler for a Specification-First Language\\[2pt]
\large Corpus-Scale Differential Measurement against Zig, GCC and Clang Pipelines}
\author{Dmitrii Vasilev\\ \small Trinity S3AI}
\date{2026-11-07}
\begin{document}
\maketitle

\begin{abstract}
t27b is the native backend of t27, a specification-first language whose source files carry their own tests. It lowers t27's internal intermediate representation to AArch64 machine code, runs a specification's tests in-process through a JIT, and emits relocatable objects and assembly listings. We measure it against the t27c reference pipeline (t27c gen, then Zig 0.16 Debug) and against t27c's C emitter compiled with GCC 12.2 and Clang 21.1, over the 1,819-specification corpus at one pinned revision. t27b passes 1,296 of the 1,311 files the reference passes (98.9\%) and agrees with the reference on all 8,344 tests both ran. Its test-compile latency, which ends with code in memory and does not link, is 4.8~ms per specification (geometric mean over 18), 8.8$\times$ below GCC -O0 compiling the generated C to an object and 640$\times$ below one Zig Debug invocation, which builds and links a test binary; when 20 specifications share one Zig invocation the gap is 8.6$\times$. The compiler is a 3.08~MB executable with one crate dependency and a median peak compile memory of 5.3~MiB. Its generated code is at the level of an unoptimising compiler: on the 20 specifications with the most GCC -O0 instructions among the 567 that every pipeline, t27b included, passes, it executes 0.91$\times$ the instructions of GCC -O0 and 3.4$\times$ those of GCC -O2, runs 4.1$\times$ slower than GCC -O2 under qemu-aarch64 emulation (3.7$\times$ slower than Clang -O2 on the 9 specifications Clang did not fold), and is 1.8$\times$ larger than GCC -O2 code; on four crypto specifications it executes 10.1$\times$ the instructions of Zig ReleaseFast. Agreement is self-consistency behind one shared front end, not independent proof: a random-program fuzzer at a nearby commit found 7 disagreements in 1,000 programs. All run times are emulated. The report and its raw JSON are committed with the repository.
\end{abstract}

\section{Introduction}

t27 is a specification-first language: a \texttt{.t27} file declares
functions and data together with the \texttt{test} and \texttt{invariant}
blocks that state what they must do. Its reference compiler, t27c, checks
a specification and emits it as Zig, C, Rust or Verilog source; the
reference way to run a specification's tests is \texttt{t27c gen} followed
by the Zig compiler. t27b is a second, native path: it mounts the same
front end and lowers the resulting internal representation directly to
AArch64 machine code, without LLVM, an assembler or a linker.

This paper asks what such a small native backend costs and what it buys
when measured over a whole corpus rather than on a chosen example. The
answer, on the evidence below, is the profile of a fast baseline compiler:
low test-compile latency and a small footprint, with generated code at the
level of an unoptimising compiler.

\paragraph{Contributions.}
(1) A corpus-scale, test-by-test differential comparison of t27b with the
reference pipeline and with C pipelines over the 1,819-specification
corpus at one revision, each pair compared on the files both ran
(Section~\ref{sec:agree}).
(2) Compile-latency, memory and footprint measurements against Zig, GCC,
Clang and rustc on the same host, including the case where Zig amortises
its fixed cost over many specifications (Sections~\ref{sec:compile}
and~\ref{sec:resources}).
(3) A code-quality measurement by executed guest instructions on
specifications ranked by GCC -O0 instruction count, a rule that uses no
t27b measurement (Section~\ref{sec:code}).
(4) An explicit list of threats to validity, including the ways in which
the agreement result is not independent (Section~\ref{sec:threats}).

\paragraph{Scope: what t27b is not.}
t27b is not an optimising compiler and not a replacement for LLVM. It is
not a multi-language backend: it does not parse C, Rust, Zig, Go or
Python, and no front end for those languages lowers to its internal
representation. In this paper ``t27 IR'' means that in-memory
representation; it is not an external interchange format, and t27b reads
no serialised IR. At the measured revision t27b did not produce a linked
executable on Linux (Section~\ref{sec:sut}).

\paragraph{Artifacts.}
Every number in this paper outside Section~\ref{sec:after} comes from
one measurement report committed to the public t27 repository,
\path{docs/reports/t27b-bench-2026-10.md}, and its raw data,
\path{docs/reports/t27b-bench-2026-10.json}
(\url{https://github.com/Trinity-S-AI-DNA/t27}). The measurements were
taken at commit \texttt{9a0f4c394} (2026-10-10; full hash in
Table~\ref{tab:env}), except two trust results that the report quotes
from the t27b lab at other commits; Section~\ref{sec:threats} names
their commits. The harness itself is lab scratch and is not committed; the
report spells out every command it ran (Section~\ref{sec:threats}). The
repository's acceptance matrix for t27b, the t27 specification
\path{specs/tri/shootout/acceptance_matrix.t27}, records the headline
ratios as cells tagged with the same commit; it is a ledger of this
report, not a second measurement.

\section{System under test}
\label{sec:sut}

\paragraph{t27c.}
Version 0.5.2, built at the measured commit. \texttt{t27c gen} emits Zig,
\texttt{gen-c} emits C and \texttt{gen-rust} emits Rust. gen-rust lowers
function bodies but no tests, so it appears only in compile-time, size and
determinism results.

\paragraph{t27b.}
t27b mounts t27c's front end (\texttt{bootstrap/src/compiler.rs}) and adds
its own AArch64 code generator. \texttt{t27b test} compiles a specification
and runs each test in an in-process JIT; with \texttt{--check} it also runs
each test in t27b's own IR interpreter and compares the two.
\texttt{--overflow trap}, the default, keeps integer-overflow and bounds
checks; \texttt{wrap} drops the overflow checks. \texttt{t27b build} writes
a relocatable object and \texttt{t27b asm} an assembly listing;
\texttt{t27b corpus} walks a directory and, with \texttt{--reference t27c},
compares every test with the reference. t27b is written in Rust. 36\% of
its backend source is generated from \texttt{.t27} specifications;
counting the front end it mounts, the generated share is 12\%. Its only
crate dependency is \texttt{serde}, used by the shared front end.

\paragraph{Limits at the measured revision.}
\texttt{t27b build} wrote Mach-O objects only, did not link, and refused
module-level variables in an object file (2 of the 20 code-quality
specifications). On aarch64-linux t27b was therefore a JIT test runner,
not a compiler that ships binaries; every t27b run time below is a JIT run.
Section~\ref{sec:after} reports what changed after the snapshot.

\section{Environment and method}

\begin{table}[h]
\centering
\small
\begin{tabularx}{\linewidth}{@{}lX@{}}
\toprule
Host & Shared x86\_64 cloud container: Intel Xeon 6975P-C, 48 host threads, cgroup quota 24 CPUs, 24 GB memory limit, Linux 6.18.5, Debian 12.15. Host 1-minute load 27--103 during the run, recorded with every sample. \\
Commit & \texttt{9a0f4c394a1474bb75103f141cc0b405d9de51b9} \\
Zig & 0.16.0; its aarch64 output goes through LLVM (LLD 21.1.0) \\
GCC & aarch64-linux-gnu-gcc 12.2.0, binutils 2.40, glibc 2.36, \texttt{-static} \\
Clang & 21.1.0 as bundled in Zig 0.16.0 (\texttt{zig cc -target aarch64-linux-musl}, \texttt{-static}) \\
rustc & 1.99.0 (LLVM 23.1.1), aarch64-unknown-linux-gnu \\
qemu & qemu-aarch64 7.2.22, user mode \\
\bottomrule
\end{tabularx}
\caption{Measurement environment.}
\label{tab:env}
\end{table}

\paragraph{Pipelines.}
t27b runs tests in its JIT under qemu-aarch64. The Zig pipelines (Debug,
ReleaseSafe, ReleaseFast) and the C pipelines (GCC and Clang at -O0 and
-O2) build static aarch64 test binaries and run each test in its own
qemu-aarch64 process, the same one-test-per-process protocol as
\texttt{t27c test-report}. Safety checks differ: t27b (trap) and Zig
Debug/ReleaseSafe check overflow and bounds at run time; Zig ReleaseFast,
GCC and Clang do not, and in C signed overflow is undefined behaviour. The
like-for-like comparison for t27b's default output is Zig ReleaseSafe.

\paragraph{Corpus and classes.}
The denominator is every \texttt{.t27} file under \texttt{specs/}: 1,819
files, none left out for being slow or broken. Each file falls into one
class per pipeline: pass (compiled, at least one test, all passed),
vacuous pass (t27b only: no runtime assert executed), compiled with no
tests, a test fails, timeout, or blocked (did not generate, did not
compile, or t27b refused a construct). Where two pipelines both ran a
file, verdicts are compared test by test, by name.

\paragraph{Compile latency.}
The host is x86\_64, so t27b was also built for the host. On x86\_64
\texttt{t27b test} runs the front end and code generation, tests included,
then stops with exit code 74 when it tries to map the JIT; the wall time of
that process is t27b's test-compile time. It therefore covers code
generation into memory, with no loading, assembling or linking, while
the Zig figures include building and linking a test binary and the C
figures stop at an object; the steps are not the same work. All other
compilers also run natively. One sample is the wall time of one process from spawn to exit,
one compile at a time, round-robin over specifications and pipelines for
6 rounds; rounds 1--5 are warm and their median is reported. Process start
(about 1~ms for t27b and t27c) is included.

\paragraph{Code-quality selection.}
Tier 1 starts from the 567 files that every pipeline, t27b included,
passes, builds each with GCC -O0, counts the guest instructions its tests
execute, and takes the 20 with the most. The ranking uses no t27b
measurement, but the pool is conditioned on t27b passing, and the rule
favours compute-heavy specifications. It selects six
benchmark kernels, three FPGA-synthesis simulations, an ISA model and ten
other specifications. gen-c output does not compile for the crypto
specifications, so they are not in tier 1; four of them are measured
against Zig only.

\paragraph{Instructions, time and size.}
A qemu TCG plugin adds each translated block's instruction count at block
entry. For Zig and C it counts between two \texttt{getppid()} markers
around each test call; for t27b it counts blocks in memory the process
made executable, which is the JIT image. The count is deterministic and
includes t27b's run-time invariants (under 6\% of its test time). Run time
is the sum over tests of \texttt{CLOCK\_MONOTONIC} brackets inside the
guest, median of 5 repetitions; these are emulated times. Code size is the
bytes of machine code for the specification's functions and tests.
Pipeline ratios are geometric means over specifications of t27b divided by
the other pipeline.

\section{Coverage and agreement}
\label{sec:agree}

\begin{table}[h]
\centering
\small
\begin{tabular}{@{}lr@{}}
\toprule
Pipeline & passes, of the 1,311 files the reference passes \\
\midrule
t27b (JIT, trap) & 1,296 (98.9\%) \\
Zig Debug, aarch64 & 1,311 (100.0\%) \\
Zig ReleaseFast, aarch64 & 1,306 (99.6\%) \\
gen-c + GCC -O0 / -O2 & 581 (44.3\%) / 578 (44.1\%) \\
gen-c + Clang -O0 / -O2 & 571 (43.6\%) / 622 (47.4\%) \\
\bottomrule
\end{tabular}
\caption{Coverage relative to the reference (t27c gen + Zig Debug, x86\_64).}
\label{tab:coverage}
\end{table}

\begin{table}[h]
\centering
\small
\begin{tabular}{@{}lrrrr@{}}
\toprule
Pair & files both ran & files that differ & tests compared & tests that differ \\
\midrule
t27b / reference & 1,315 & 0 & 8,344 & 0 \\
t27b / Zig ReleaseFast & 1,315 & 5 & 8,344 & 8 \\
Zig Debug / Zig ReleaseFast & 1,337 & 7 & 8,531 & 11 \\
GCC -O0 / GCC -O2 & 644 & 4 & 4,906 & 4 \\
Clang -O0 / Clang -O2 & 635 & 39 & 4,798 & 107 \\
GCC -O2 / Clang -O2 & 635 & 35 & 4,798 & 95 \\
t27b / GCC -O2 & 610 & 24 & 3,790 & 48 \\
t27b / Clang -O2 & 634 & 12 & 3,901 & 28 \\
\bottomrule
\end{tabular}
\caption{Test-by-test disagreements.}
\label{tab:agree}
\end{table}

t27b agrees with the reference on every test both ran: 0 of 8,344
(Table~\ref{tab:agree}). The reference itself is unchanged by moving from
x86\_64 to aarch64 (0 of 8,513 tests differ). Zig ReleaseFast is not the
reference: it flips 7 files against Zig Debug; t27b runs 5 of them and
agrees with Debug on all 5. The C pipelines disagree with themselves
across optimisation levels, because gen-c output contains undefined
behaviour; GCC -O2 and Clang -O2 differ from each other on 35 files. So
t27b's differences with the optimised C pipelines (24 files against GCC
-O2, 12 against Clang -O2) do not by themselves locate a t27b fault. The
C rows cover only the files gen-c compiles, about a third of the corpus.

Of the 15 files the reference passes and t27b does not, 14 are blocked
(i128 three times, an escaping frame address three times, an empty typed
array literal twice, \texttt{@cos}, \texttt{anytype}, \texttt{std.*}, an
assignment form, a type mismatch, and a 1~MiB frame-size limit) and 1
passes vacuously. t27b passes 3 files the reference cannot compile.
A native arm64 CI run at the same commit printed the same t27b tally as
the emulated run.

\section{Compile latency}
\label{sec:compile}

\begin{table}[h]
\centering
\small
\begin{tabular}{@{}lrr@{}}
\toprule
Pipeline (18 tier-1 specifications) & geomean warm median (ms) & $\times$ t27b test-compile \\
\midrule
t27b test-compile (tests included) & 4.75 & 1.00 \\
t27b build (object, no tests) & 5.11 & 1.08 \\
gen-c + GCC -O0, object & 42.0 & 8.84 \\
gen-c + GCC -O2, object & 100 & 21.0 \\
gen-c + Clang -O0, object & 71.1 & 15.0 \\
gen-c + Clang -O2, object & 105 & 22.1 \\
gen + Zig Debug, test binary & 3,050 & 642 \\
gen + Zig ReleaseFast, test binary & 13,808 & 2,906 \\
\bottomrule
\end{tabular}
\caption{Compile latency per specification, native x86\_64. The 18
specifications are the tier-1 set minus the two \texttt{t27b build}
refuses; generator time is included where there is a generator.}
\label{tab:compile}
\end{table}

Most of Zig's time per specification is fixed cost (test runner, standard
library machinery and the LLVM link), which a project amortises by
compiling many specifications at once. With all 20 tier-1 specifications
in one \texttt{zig test} invocation, Zig takes 6.21~s in Debug and 32.9~s
in ReleaseFast, against 0.72~s for 20 separate t27b test-compile
processes: the gap falls from about 640$\times$ to 8.6$\times$ (Debug) and
from about 2,900$\times$ to 46$\times$ (ReleaseFast). t27b has no
multi-file mode. Over the 567 files every pipeline passes, t27b
test-compiles 183 files per second (3.10~s in all), against 28 files per
second for gen-c + GCC -O0 to objects (6.4$\times$ slower). The first Zig
build after an empty global cache costs more again (6.2~s in Debug on one
kernel); t27b, GCC and rustc have no such step.

\section{Generated code}
\label{sec:code}

\begin{table}[h]
\centering
\small
\begin{tabular}{@{}lrrrr@{}}
\toprule
t27b (trap) / pipeline & $n$ & executed instructions & emulated run time & code size \\
\midrule
Zig Debug & 20 & 0.54 & 0.69 & 0.66 \\
Zig ReleaseSafe & 20 & 3.17 & 3.41 & 1.89 \\
Zig ReleaseFast & 20 & 4.47 & 4.14 & 2.03 \\
GCC -O0 & 20 & 0.91 & 1.11 & 1.34 \\
GCC -O2 & 20 & 3.43 & 4.05 & 1.82 \\
Clang -O0 & 20 & 0.41 & 0.62 & 0.65 \\
Clang -O2 & 9 & 3.50 & 3.72 & 1.69 \\
\bottomrule
\end{tabular}
\caption{Geometric-mean ratios over tier 1. Below 1, t27b is better.
Clang -O2 evaluated the test at compile time on 11 of the 20
specifications (it executed under 1\% of GCC -O0's instructions); those
are left out of its row.}
\label{tab:code}
\end{table}

t27b's code is at the level of an unoptimising compiler
(Table~\ref{tab:code}). It executes about as many instructions as GCC -O0
and beats only unoptimised LLVM output (Zig Debug, Clang -O0). On tier 1,
every optimising pipeline is 3 to 5 times better by geometric mean
(individual specifications vary more): against the like-for-like
checked build, Zig ReleaseSafe, t27b executes 3.2$\times$ the instructions
and takes 3.4$\times$ the emulated time; against GCC -O2, 3.4$\times$ and
4.1$\times$; against Clang -O2 where it did not fold the test,
3.5$\times$ and 3.7$\times$. Its code is larger than every optimising
build. The gap is larger outside tier 1: on the four crypto
specifications measured against Zig only, t27b executes 7.0$\times$ the
instructions of Zig ReleaseSafe and 10.1$\times$ those of Zig ReleaseFast.

Overflow checks add about 12\% to the instructions of t27b's unchecked
build (trap/wrap 1.12$\times$) and 24\% to its emulated time. One case is pathological: before every test the JIT copies
the whole module-variable image so each test starts from fresh state, and
a specification with large module arrays (\texttt{compiler/core/t27core},
95 tests) runs 8.3$\times$ slower than Zig Debug although t27b executes
fewer instructions. Instruction counts are the better measure here:
emulation cost varies by instruction mix, and on one kernel Zig
ReleaseFast executes 43\% fewer instructions than t27b yet takes 25\%
longer under qemu.

\section{Compiler resources}
\label{sec:resources}

\begin{table}[h]
\centering
\small
\begin{tabular}{@{}lrrr@{}}
\toprule
Compiler step & $n$ & median peak RSS (MiB) & max (MiB) \\
\midrule
t27b test-compile & 24 & 5.26 & 48.4 \\
t27b build & 22 & 5.12 & 14.1 \\
Zig Debug / ReleaseFast & 24 & 354 / 500 & 412 / 670 \\
GCC -O0 / -O2, object & 20 & 27.1 / 31.6 & 59.4 / 173 \\
Clang -O0 / -O2, object & 20 & 68.0 / 75.5 & 118 / 243 \\
rustc -O0 / -O3, object & 18 & 112 / 120 & 138 / 134 \\
\bottomrule
\end{tabular}
\caption{Peak resident memory of each compiler's whole process tree, from
\texttt{wait4}. $n$ is the number of code-quality specifications (tier 1
plus tier 2, up to 24) that step compiled; failed compiles are left out,
so the rows are not over the same set. Compile steps only: running the
tests, aarch64 t27b under qemu peaks at a median of 26.1~MiB (max 119~MiB,
qemu's own memory included).}
\label{tab:mem}
\end{table}

t27b is one executable: 3.08~MB for x86\_64 and 2.79~MB for aarch64,
unstripped, with Rust's standard library linked statically and no
assembler or linker. For comparison, the t27c executable is 11.9~MB, Zig
0.16.0 (which also provides the Clang used here) is 363~MB installed, GCC
12 for aarch64 is about 145~MB, and rustc 1.99 is 790~MB. The comparison
is not like for like: those toolchains produce linked executables, which
t27b did not at the measured revision.

\section{Determinism}
\label{sec:determinism}

Each tier-1 specification was built by every pipeline three times, twice
in one directory and once from a different absolute path, where the
pipeline could build it: \texttt{t27b build} and \texttt{t27b asm} built 18
of the 20 (they refuse the two with module-level variables) and rustc 15.
All outputs that were built were bit-identical within one directory. t27b's objects and assembly
listings, t27c's generators, GCC (without \texttt{-g}) and rustc were also
identical across directories; Zig and \texttt{zig cc} record the build
directory in DWARF and were not.

\section{Threats to validity}
\label{sec:threats}

\paragraph{Self-consistency, not independence.}
The reference t27b agrees with is produced by the same t27c front end
that t27b mounts, and the JIT/interpreter check reads one IR. A
diverse-double-compiling receipt over three routes (lab, commit
\texttt{0e8564505}) ends in backend agreement, not independence, for the
same reason~\cite{wheeler}. A source-level random-program fuzzer against
the reference reaches beyond the corpus and, in the lab run the report
quotes (commit \texttt{8e2e44b7e}, 1,000 programs, 1,975 tests), found 7
disagreements; the corpus result of Table~\ref{tab:agree} therefore does
not mean t27b agrees with the reference on every program. No check here compares t27b with an independently written
t27 front end.

\paragraph{Emulation and a shared host.}
All run times are qemu-aarch64 on x86\_64; they rank pipelines on this
host and are not arm64 speed. The host was shared with other jobs (load
27--103 on 48 threads); medians are reported, but millisecond-scale
compile samples are noisy (over 5 runs, the median max/min spread was
34\% for t27b and 36\% for GCC -O0).

\paragraph{Different contracts.}
t27b re-initialises module variables before each test and runs invariants
at run time; the Zig and C runners do neither. Undefined behaviour in
gen-c output changes C verdicts with the optimisation level, and an
optimiser may exploit it. Clang is Zig's bundled Clang with musl, not a
distribution Clang with glibc.

\paragraph{Which Zig backend.}
Zig 0.15 made its self-hosted x86-64 backend the default for Debug builds,
and its release notes describe Debug compilation as five times faster with
it~\cite{zig015}. The aarch64 target measured here goes through LLVM
(Table~\ref{tab:env}), so the 640$\times$ figure is against Zig Debug
through LLVM and does not carry over to Zig's self-hosted backends.

\paragraph{Selection and repetition.}
Tier 1 excludes everything gen-c cannot compile. Coverage numbers come
from one corpus run per pipeline. The 4.8~ms figure is a geometric mean
over 18 specifications, not over the corpus.

\paragraph{Reproducibility.}
The harness (about 1,650 lines of lab scratch) is not committed; the
report lists every command it ran and the tool hashes, and the JSON holds
every sample with its host load. An earlier single-specification
demonstration on an arm64 workstation is not used here, because its
specification is not on the main branch and no raw runs were retained.

\section{State after the snapshot}
\label{sec:after}

The results above refer to commit \texttt{9a0f4c394}. The following
changes are on the main branch, are reported in the repository's issues,
and were not re-measured for this paper. \texttt{t27b build --elf} writes
ELF64 relocatable objects for aarch64-linux, which supersedes the
Mach-O-only limit of Section~\ref{sec:sut}; \texttt{t27b run} executes a
named entry point; \texttt{t27b verilog} emits RTL; scalar exports follow
the AAPCS64 C calling convention. The corpus has also grown. Later lab
runs at other commits are not used here: the revisions are not
interchangeable, and no post-snapshot figure enters this paper.

\section{Claims and evidence}

\begin{table}[h]
\centering
\small
\begin{tabularx}{\linewidth}{@{}XlX@{}}
\toprule
Claim & Report section & Status \\
\midrule
Coverage 1,296 / 1,311; 0 of 8,344 tests disagree & A & measured; self-consistency \\
Test-compile 4.8~ms; 8.8$\times$ GCC -O0; 640$\times$, batched 8.6$\times$, Zig Debug & B & measured; shared host \\
0.91$\times$ GCC -O0 and 3.4$\times$ GCC -O2 instructions; 4.1$\times$ / 3.7$\times$ emulated time vs GCC / Clang -O2 & C & measured; emulated time \\
5.3~MiB median peak memory; 3.08~MB executable & D & measured \\
Bit-identical, path-independent t27b outputs & E & measured \\
Native run time comparable to Clang -O2 & -- & not claimed; contradicted by C \\
External IR input; C, Rust, Zig, Go or Python front ends & -- & not implemented \\
\bottomrule
\end{tabularx}
\caption{Claims and their evidence in \texttt{docs/reports/t27b-bench-2026-10.md}.}
\label{tab:claims}
\end{table}

\section{Related work}

LLVM established the pattern of shared front ends, a common IR and
retargetable optimising backends~\cite{lattner2004}; t27b is the opposite
trade, a single-target baseline backend with no optimiser. Cranelift,
offered as an alternative code generator for rustc, makes a similar
compile-time versus code-quality trade within an existing
compiler~\cite{rustcCodegenSsa,craneliftTradeoff}.

\paragraph{Baseline compilers.}
Fast non-optimising compilers are an established design point. TinyCC
reports compiling, assembling and linking C several times faster than
GCC~\cite{tcc}; QBE aims at 70\% of the performance of industrial
optimising compilers in 10\% of the code~\cite{qbe}. In WebAssembly
engines, V8's Liftoff makes a single pass over a function's
bytes~\cite{liftoff}, and Wasmtime's Winch names compilation performance
as its primary goal, with a single pass and no intermediate
representation~\cite{winch}. Titzer compares six baseline compilers and
the tradeoff between compilation speed and code
quality~\cite{titzer2023}. Copy-and-patch compilation reports generating
code 4.9--6.5$\times$ faster than Liftoff~\cite{copypatch}; TPDE, a
back-end framework for LLVM IR, reports compiling 8--24$\times$ faster
than LLVM -O0 with on-par run time on SPECint 2017~\cite{tpde}; Umbra's
Flying Start back end targets low query latency~\cite{umbra}. t27b differs from
these in what it compiles and checks: it lowers one specification
language's own IR and runs that language's embedded tests in process. None
of these systems was measured here, and the ratios in this paper are
against the same specifications' C and Zig outputs only.

\paragraph{Testing compilers.}
Csmith and EMI find
compiler bugs by random programs and equivalence modulo
inputs~\cite{csmith,emi}; this paper runs neither, and its differential
method compares two backends behind one front end. The corpus also
contains specifications of number formats described
elsewhere~\cite{goldenfloat,goldenruler}; two of them (GF16 and OCP MX)
are among the measured specifications.

\section{Future work}

The measured gaps name the next steps: middle-end optimisation passes to
close the 3--5$\times$ tier-1 instruction gap (10$\times$ on the crypto
specifications); a multi-file build mode to keep
the latency lead when Zig amortises; per-test state isolation without
copying the whole module image; a differential fuzzing gate with zero
disagreements; and native arm64 timings in place of emulated ones.
Lowering restricted subsets of other languages through a versioned IR
contract to t27b is an architectural goal; no such front end exists.

\section{Conclusion}

Over the full 1,819-specification corpus at one revision, t27b passes
98.9\% of what its reference passes and agrees with it on all 8,344 tests
both ran. It test-compiles a specification in about 5~ms with about
5~MiB of memory from a 3~MB executable without LLVM, though without
linking; its code executes 3--5$\times$ the instructions of optimising
compilers on tier 1 by geometric mean, and 10$\times$ on the crypto
specifications. That is the profile
of a fast baseline compiler; its agreement with the reference is
self-consistency behind a shared front end, not independent verification.

\begin{thebibliography}{99}

\bibitem{goldenfloat}
D.~Vasilev,
\emph{GoldenFloat: A Phi-Derived Static-Split Floating-Point Family from GF4 to GF1024 with a Lucas-Exact Integer Identity},
2026.
\href{https://arxiv.org/abs/2606.05017}{arXiv:2606.05017}.

\bibitem{goldenruler}
D.~Vasilev,
\emph{Golden Ruler: A Numeric Format Catalog with Bit-Exact Conformance Vectors for FP8, BF16, MXFP4, and Microscaling Formats},
2026.
\href{https://arxiv.org/abs/2606.09686}{arXiv:2606.09686}.

\bibitem{lattner2004}
C.~Lattner and V.~Adve,
\emph{LLVM: A Compilation Framework for Lifelong Program Analysis \& Transformation},
CGO 2004.
DOI 10.1109/CGO.2004.1281665.

\bibitem{rustcCodegenSsa}
Rust Compiler Development Guide,
\emph{Backend-agnostic code generation}.
\url{https://rustc-dev-guide.rust-lang.org/backend/backend-agnostic.html}.

\bibitem{craneliftTradeoff}
\emph{Using rustc\_codegen\_cranelift},
Inside Rust blog, 2020-11-15.
\url{https://blog.rust-lang.org/inside-rust/2020/11/15/Using-rustc_codegen_cranelift/}.

\bibitem{csmith}
X.~Yang, Y.~Chen, E.~Eide, and J.~Regehr,
\emph{Finding and Understanding Bugs in C Compilers},
PLDI 2011.
DOI 10.1145/1993498.1993532.

\bibitem{emi}
V.~Le, M.~Afshari, and Z.~Su,
\emph{Compiler Validation via Equivalence Modulo Inputs},
PLDI 2014.
DOI 10.1145/2594291.2594334.

\bibitem{wheeler}
D.~A.~Wheeler,
\emph{Fully Countering Trusting Trust through Diverse Double-Compiling},
PhD dissertation, George Mason University, 2009.
\href{https://arxiv.org/abs/1004.5534}{arXiv:1004.5534}.

\bibitem{tcc}
F.~Bellard,
\emph{Tiny C Compiler}.
\url{https://bellard.org/tcc/}.

\bibitem{qbe}
\emph{QBE: compiler backend}.
\url{https://c9x.me/compile/}.

\bibitem{liftoff}
C.~Backes,
\emph{Liftoff: a new baseline compiler for WebAssembly in V8},
V8 blog, 2018-08-20.
\url{https://v8.dev/blog/liftoff}.

\bibitem{winch}
Bytecode Alliance,
\emph{Winch: WebAssembly Intentionally Non-optimizing Compiler and Host},
Wasmtime repository.
\url{https://github.com/bytecodealliance/wasmtime/tree/main/winch}.

\bibitem{titzer2023}
B.~L.~Titzer,
\emph{Whose baseline compiler is it anyway?},
2023.
\href{https://arxiv.org/abs/2305.13241}{arXiv:2305.13241}.

\bibitem{copypatch}
H.~Xu and F.~Kjolstad,
\emph{Copy-and-Patch Compilation: A fast compilation algorithm for high-level languages and bytecode},
2021.
DOI 10.1145/3485513;
\href{https://arxiv.org/abs/2011.13127}{arXiv:2011.13127}.

\bibitem{tpde}
T.~Schwarz, T.~Kamm, and A.~Engelke,
\emph{TPDE: A Fast Adaptable Compiler Back-End Framework},
2025.
\href{https://arxiv.org/abs/2505.22610}{arXiv:2505.22610}.

\bibitem{umbra}
T.~Kersten, V.~Leis, et al.,
\emph{Tidy Tuples and Flying Start: fast compilation and fast execution of relational queries in Umbra},
The VLDB Journal 30:883--905, 2021.
DOI 10.1007/s00778-020-00643-4.

\bibitem{zig015}
\emph{Zig 0.15.1 Release Notes}, section ``x86 Backend''.
\url{https://ziglang.org/download/0.15.1/release-notes.html}.

\end{thebibliography}

\end{document}
